\chapter{The Prime Number Theorem}

\begin{center}\itshape
I soon recognized that behind all of its fluctuations, this frequency is on the average inversely proportional to the logarithm.\\
\upshape --- Gauss, letter to Encke (1849), on the tables of primes he had made as a boy
\end{center}

Some people do crosswords. Gauss, in idle quarters of an hour, counted primes: by the time he wrote the letter above he had tabulated them up to three
million. Around 1792, at fifteen, he counted the primes in blocks of a thousand and guessed that near $x$ they have density about $\frac1{\log x}$, so that
\[
\pi(x)\approx\operatorname{Li}(x)=\int_2^x\frac{\mathrm dt}{\log t}.
\]
Legendre proposed $\frac x{\log x-1.08366}$ a few years later. Chebyshev proved in 1852 that $\frac{\pi(x)\log x}x$ stays between two constants close to $1$ and that its
limit, if it has one, is $1$ (Corollary~\ref{cor:chebratio}). Riemann's memoir of 1859 tied the question to the zeros of $\zeta$, and in 1896 Hadamard and de la Vallée Poussin,
independently, proved the \emph{prime number theorem}:
\[
\pi(x)\sim\frac x{\log x}.
\]
The key, in both proofs, is that $\zeta$ has no zero on the line $\operatorname{Re}s=1$. This chapter proves that fact and then deduces the theorem by D.~J.~Newman's
argument of 1980, in the form given by Zagier in 1997: a Tauberian theorem with a three-line contour integral at its heart. The same method gives, at no extra cost,
$\sum\frac{\mu(n)}n=0$, the theorem for arithmetic progressions, and the limits that Chapters~1 and~4 left open.

\begin{think}
At $x=10^8$, $\frac x{\log x}$ is too small by $6\%$, while $\operatorname{Li}(x)$ is too large by about $0.013\%$ (Figure~\ref{fig:pnt}). Both are asymptotic to $\pi(x)$, if one is.
Why might $\operatorname{Li}$ be the ``right'' main term? Compare $\operatorname{Li}(x)$ with $\frac x{\log x}$ by integrating by parts.
\end{think}

\begin{figure}[ht]
\centering
\fig{ch05-pnt}
\caption{$\frac{\pi(x)\log x}x$ (coral) and $\frac{\pi(x)}{\operatorname{Li}(x)}$ (charcoal) for $10^2\le x\le10^8$, in logarithmic scale. Both tend to $1$; the first very slowly, from above
($1.0613$ at $10^8$), the second fast, from below ($0.99987$ at $10^8$, where $\pi(x)=5\,761\,455$ and $\operatorname{Li}(x)-\pi(x)=753$).}
\label{fig:pnt}
\end{figure}

\section{The logarithmic derivative}

As in Chapter~1, $\psi(x)=\sum_{n\le x}\Lambda(n)$ and $\theta(x)=\sum_{p\le x}\log p$. The analytic object attached to $\psi$ is the logarithmic derivative of $\zeta$. We treat
$L$-functions at the same time, since it costs nothing; with $q=1$ and $\chi=1$ the reader may read $\zeta$ throughout.

\begin{lemma}\label{lem:logderiv}
Let $\chi$ be a character modulo $q$. For $\operatorname{Re}s>1$,
\[
-\frac{L'(s,\chi)}{L(s,\chi)}=\sum_{n\ge1}\frac{\chi(n)\Lambda(n)}{n^s}.
\]
In particular $-\frac{\zeta'(s)}{\zeta(s)}=\sum_n\Lambda(n)n^{-s}$.
\end{lemma}
\begin{proof}
By Theorem~\ref{thm:lbasic}(a), $L(s,\chi)=\exp(E(s))$ with $E(s)=\sum_p\sum_{k\ge1}\frac{\chi(p)^k}{kp^{ks}}$. On $\operatorname{Re}s\ge\sigma_0>1$ the terms, which are entire functions, are
bounded by $p^{-k\sigma_0}$, whose sum is finite, so the series converges uniformly there and may be differentiated term by term (Volume~5, Theorem~11.15(a)):
$E'(s)=-\sum_p\sum_k\chi(p^k)(\log p)p^{-ks}$, an absolutely convergent series. Grouping its terms by $n=p^k$ gives $-\sum_n\chi(n)\Lambda(n)n^{-s}$. Finally
$L'=E'e^E=E'L$.
\end{proof}

\section{No zeros on the line \texorpdfstring{$\operatorname{Re}s=1$}{Re s = 1}}

Recall that $L(s,\chi)$ is holomorphic on $\operatorname{Re}s>0$ for $\chi\ne\chi_0$, and that $L(s,\chi_0)$, in particular $\zeta(s)$, is holomorphic there except for a simple pole at $1$
(Theorem~\ref{thm:lbasic}). So $L(1+it,\chi)$ makes sense except for $\chi=\chi_0$ and $t=0$.

\begin{theorem}[Hadamard, de la Vallée Poussin, 1896]\label{thm:line}
Let $\chi$ be a character modulo $q$ and $t$ real, with $(\chi,t)\ne(\chi_0,0)$. Then $L(1+it,\chi)\ne0$. In particular $\zeta(1+it)\ne0$ for every real $t\ne0$.
\end{theorem}
\begin{proof}
\emph{The inequality.} For every real $\phi$,
\[
\begin{aligned}3+4\cos\phi+\cos2\phi&=3+4\cos\phi+2\cos^2\phi-1\eqbr=2(1+\cos\phi)^2\ge0.\end{aligned}
\]
Let $\sigma>1$. Since $|e^w|=e^{\operatorname{Re}w}$, Theorem~\ref{thm:lbasic}(a) gives $\log|L(s,\chi)|=\operatorname{Re}E(s)$ with the notation of the last proof. Apply it to $\chi_0$ at $\sigma$, to $\chi$
at $\sigma+it$ and to $\chi^2$ at $\sigma+2it$, and set $P(\sigma)=|L(\sigma,\chi_0)|^3|L(\sigma+it,\chi)|^4|L(\sigma+2it,\chi^2)|$. For $p\nmid q$ write $\chi(p)^kp^{-ikt}=e^{i\phi}$ with
$\phi=\phi_{p,k}$ real; the terms with $p\mid q$ vanish. Then
\[
\log P(\sigma)=\sum_{p\nmid q}\sum_{k\ge1}\frac{3+4\cos\phi_{p,k}+\cos2\phi_{p,k}}{kp^{k\sigma}}\ge0,
\]
that is, $P(\sigma)\ge1$ for every $\sigma>1$.

\emph{The contradiction.} The case $t=0$ is Theorem~\ref{thm:lnonzero}, so let $t\ne0$ and suppose that $L(1+it,\chi)=0$. Let $\sigma\to1^+$.
\begin{itemize}
\item $0<L(\sigma,\chi_0)\le\zeta(\sigma)\le\frac\sigma{\sigma-1}$, as in the proof of Theorem~\ref{thm:lnonzero}, so $L(\sigma,\chi_0)^3=O\bigl((\sigma-1)^{-3}\bigr)$.
\item $L(s,\chi)$ is holomorphic near $1+it$ (for $\chi=\chi_0$ because $t\ne0$) and vanishes there, so $|L(\sigma+it,\chi)|\le C(\sigma-1)$ and its fourth power is $O\bigl((\sigma-1)^4\bigr)$.
\item $L(s,\chi^2)$ is holomorphic near $1+2it$, since $2t\ne0$, so it is bounded there.
\end{itemize}
Hence $P(\sigma)=O(\sigma-1)$ tends to $0$, which contradicts $P(\sigma)\ge1$.
\end{proof}

\begin{think}
The inequality $1+\cos\phi\ge0$ gives $|\zeta(\sigma)\zeta(\sigma+it)|\ge1$ in the same way. Why is that useless? What do the exponents $3$ and $4$ achieve, and would $3+4\cos\phi+\cos2\phi$
work if $\zeta$ had a double pole at $1$?
\end{think}

\begin{remark}
The weights $3,4,1$ were found by Mertens in 1898, when he simplified de la Vallée Poussin's proof. The argument shows more than non-vanishing: it gives a quantitative
lower bound for $|\zeta(1+it)|$, and with Jensen's formula it pushes the zeros a little to the left of the line, into a region $\sigma>1-\frac c{\log|t|}$. Chapter~6 proves this
\emph{zero-free region}, and Chapter~7 turns it into an error term for the prime number theorem.
\end{remark}

\section{Newman's Tauberian theorem}

This section is the engine of the chapter. For eighty years, proofs of the prime number theorem filled a chapter of hard analysis; in 1980 Donald Newman
found a contour integral that does the job in a page, and Zagier polished it into the form below. A \emph{Tauberian} theorem goes from the behaviour of a transform back to the behaviour of a function, which needs some extra hypothesis (here: boundedness).
The function below is a \emph{Laplace transform}. We say that $f\colon[0,\infty[\to\C$ is \emph{piecewise continuous} if on every bounded interval it is continuous except at finitely
many points, where it has left and right limits.

\begin{theorem}[Newman, 1980]\label{thm:newman}
Let $f\colon[0,\infty[\to\C$ be bounded and piecewise continuous, and for $\operatorname{Re}z>0$ let
\[
G(z)=\int_0^\infty f(t)e^{-zt}\dd t.
\]
Suppose that there is a function $g$, holomorphic on an open set $U$ containing the closed half-plane $\operatorname{Re}z\ge0$, with $g=G$ on $\operatorname{Re}z>0$. Then the integral
$\int_0^\infty f(t)\dd t$ converges, and equals $g(0)$.
\end{theorem}

The hypothesis is about $G$ on the imaginary axis, where the integral defining it need not converge; the conclusion is the convergence of that integral at the single point $z=0$.

\begin{proof}
Let $B=\sup|f|$. For $\operatorname{Re}z=x>0$ the integral defining $G$ converges absolutely, as $|f(t)e^{-zt}|\le Be^{-xt}$. For $T>0$ let
\[
g_T(z)=\int_0^Tf(t)e^{-zt}\dd t.
\]
Split $[0,T]$ at the finitely many discontinuities of $f$: on each piece $f$ extends to a continuous function on the closed piece, and the integrand is bounded by
$Be^{T|z|}$, so Volume~5, Theorem~11.15(b) shows that each piece, hence $g_T$, is an entire function. We must show that $g_T(0)=\int_0^Tf\to g(0)$ as $T\to\infty$.

\emph{The contour.} Fix $R>0$. The closed half-disc $\{|z|\le R,\operatorname{Re}z\ge0\}$ is compact and contained in $U$, so it is at a positive distance $\eta$ from the closed set
$\C\setminus U$. Choose $0<\delta<\min(\eta,R)$. Every $z=x+iy$ with $|z|\le R$ and $-\delta\le x<0$ is at distance $|x|<\eta$ from the point $iy$ of the half-disc, so the closure of
\[
D=\{z:|z|<R,\ \operatorname{Re}z>-\delta\}
\]
lies in $U$. $D$ is convex; let $C$ be its positive boundary. Split $C$ into $C_+$, the half-circle $|z|=R$, $\operatorname{Re}z\ge0$, and $C_-$, the rest, which lies in $\operatorname{Re}z<0$
except for its endpoints $\pm iR$ (Figure~\ref{fig:contour}). Let $C_-'$ be the half-circle $|z|=R$, $\operatorname{Re}z\le0$, run from $iR$ to $-iR$ like $C_-$. Put
\[
K(z)=e^{zT}\Bigl(1+\frac{z^2}{R^2}\Bigr)\frac1z.
\]
The function $g(z)K(z)$ is holomorphic on $U\setminus\{0\}$ with a simple pole at $0$ of residue $g(0)$, and $g_T(z)K(z)$ is holomorphic on $\C\setminus\{0\}$ with residue $g_T(0)$. The
residue theorem (Volume~5, Theorem~11.26), applied to $D$ and to the disc $|z|<R$, gives
\[
\begin{aligned}2\pi i\bigl(g(0)-g_T(0)\bigr)&=\int_CgK-\int_{|z|=R}g_TK\\&=\int_{C_+}(g-g_T)K+\int_{C_-}gK\eqbr\quad-\int_{C_-'}g_TK.\end{aligned}
\]

\emph{The factor $1+\frac{z^2}{R^2}$.} On $|z|=R$ we have $R^2=z\bar z$, so $1+\frac{z^2}{R^2}=\frac{z(\bar z+z)}{R^2}=\frac{2xz}{R^2}$ with $x=\operatorname{Re}z$, and therefore
\[
|K(z)|=e^{xT}\frac{2|x|}{R^2}\qquad(|z|=R).
\]
This is the whole trick: the factor vanishes on the imaginary axis exactly where the other bounds blow up.

\emph{On $C_+$.} For $x>0$, $|g(z)-g_T(z)|=\bigl|\int_T^\infty f(t)e^{-zt}\dd t\bigr|\le B\int_T^\infty e^{-xt}\dd t=\frac{Be^{-xT}}x$. So $|(g-g_T)K|\le\frac{2B}{R^2}$ on $C_+$ (at the two endpoints too,
by continuity), and the integral over $C_+$, of length $\pi R$, is at most $\frac{2\pi B}R$.

\emph{On $C_-'$.} For $x<0$, $|g_T(z)|\le B\int_0^Te^{-xt}\dd t\le B\int_{-\infty}^Te^{-xt}\dd t=\frac{Be^{-xT}}{|x|}$. So again $|g_TK|\le\frac{2B}{R^2}$, and the integral over $C_-'$ is at most $\frac{2\pi B}R$.

\emph{On $C_-$.} Here $g$ is no longer compared with $g_T$, and the bound must come from $e^{zT}$. The function $h(z)=g(z)(1+\frac{z^2}{R^2})\frac1z$ does not depend on $T$ and is
continuous on $C_-$ (which lies in $U$, at distance at least $\delta$ from $0$), so $|h|\le H$ there; and $|gK|=|h|e^{xT}$. Let $0<\varepsilon<\delta$. The part of $C_-$ where
$x>-\varepsilon$ consists of two small arcs near $\pm iR$, of total length $\ell(\varepsilon)$, which tends to $0$ with $\varepsilon$; on the rest, $|gK|\le He^{-\varepsilon T}$. With $\lambda$ the length
of $C_-$,
\[
\Bigl|\int_{C_-}gK\Bigr|\le H\ell(\varepsilon)+H\lambda e^{-\varepsilon T},
\]
so $\limsup_{T\to\infty}\bigl|\int_{C_-}gK\bigr|\le H\ell(\varepsilon)$ for every $\varepsilon$, and the limit is $0$.

\emph{Conclusion.} $\limsup_{T\to\infty}|g(0)-g_T(0)|\le\frac1{2\pi}\cdot\frac{4\pi B}R=\frac{2B}R$. As $R$ was arbitrary, $g_T(0)\to g(0)$.
\end{proof}

\begin{figure}[ht]
\centering
\fig{ch05-contour}
\caption{Newman's contour. $C$ bounds the shaded region $|z|<R$, $\operatorname{Re}z>-\delta$, in which $g$ is holomorphic. On the right half-circle $C_+$, $g-g_T$ is a tail of the
integral and is small. On the left, $g_T$ is integrated over the half-circle $C_-'$ (dashed), and $g$ over $C_-$, where $e^{zT}$ kills everything as $T\to\infty$.}
\label{fig:contour}
\end{figure}

\begin{beauty}
Nothing in the proof knows about primes. The arithmetic enters only through the hypothesis: $G$ continues across the imaginary axis. For $\zeta$ this will be the absence of zeros
on $\operatorname{Re}s=1$. The deep fact and the soft analysis are separated as cleanly as they can be.
\end{beauty}

\section{The prime number theorem}

All the pieces are on the table; it remains to put them together. To go from the convergence of an integral to an asymptotic, we use that $\psi$ increases.

\begin{lemma}\label{lem:monotone}
Let $A\colon[1,\infty[\to\R$ be non-decreasing and $c>0$. If the integral $\displaystyle\int_1^\infty\frac{A(x)-cx}{x^2}\dd x$ converges, then $A(x)\sim cx$.
\end{lemma}
\begin{proof}
Let $\lambda>1$ and suppose that $A(x)\ge\lambda cx$ for arbitrarily large $x$. For such an $x$ and $t\in[x,\lambda x]$, $A(t)\ge A(x)\ge\lambda cx$, so, with $t=ux$,
\[
\begin{aligned}\int_x^{\lambda x}\frac{A(t)-ct}{t^2}\dd t&\ge c\int_x^{\lambda x}\frac{\lambda x-t}{t^2}\dd t\eqbr=c\int_1^\lambda\frac{\lambda-u}{u^2}\dd u>0.\end{aligned}
\]
The right side does not depend on $x$, while the left side tends to $0$ as $x\to\infty$ by the Cauchy criterion for the convergent integral: a contradiction. So
$\limsup\frac{A(x)}{cx}\le\lambda$. Similarly, if $\lambda<1$ and $A(x)\le\lambda cx$ for arbitrarily large $x$, then $A(t)\le\lambda cx$ on $[\lambda x,x]$ and
$\int_{\lambda x}^x\frac{A(t)-ct}{t^2}\dd t\le c\int_\lambda^1\frac{\lambda-u}{u^2}\dd u<0$, again a contradiction; so $\liminf\frac{A(x)}{cx}\ge\lambda$. Let $\lambda\to1$.
\end{proof}

Write $Z(s)=(s-1)\zeta(s)$. By Volume~5, Proposition~12.6, $Z(s)=1+(s-1)\bigl(\zeta(s)-\frac1{s-1}\bigr)$ is holomorphic on $\operatorname{Re}s>0$, with $Z(1)=1$ and $Z'(1)=\gamma$. By
Theorem~\ref{thm:line} and the Euler product, $Z(s)\ne0$ on the closed half-plane $\operatorname{Re}s\ge1$; as $Z$ is continuous, the set $V=\{\operatorname{Re}s>0:Z(s)\ne0\}$ is open and contains that half-plane.

\begin{theorem}[the prime number theorem]\label{thm:pnt}
As $x\to\infty$,
\[
\begin{gathered}\psi(x)\sim x,\qq\theta(x)\sim x,\qq\pi(x)\sim\frac x{\log x}.\end{gathered}
\]
Moreover the integral $\displaystyle\int_1^\infty\frac{\psi(x)-x}{x^2}\dd x$ converges, and equals $-1-\gamma$.
\end{theorem}
\begin{proof}
Let $f(t)=\psi(e^t)e^{-t}-1$ for $t\ge0$. It is piecewise continuous ($\psi$ is a step function, with finitely many jumps on each bounded interval) and bounded, since
$\psi(x)=O(x)$ (Proposition~\ref{prop:chebyshev}). With $x=e^t$, $\int_0^Tf(t)\dd t=\int_1^{e^T}\frac{\psi(x)-x}{x^2}\dd x$, so we must show that Newman's theorem applies, with
$g(0)=-1-\gamma$.

\emph{The transform.} Let $\operatorname{Re}z>0$ and $s=z+1$. Then $G(z)=\int_1^\infty(\psi(x)-x)x^{-s-1}\dd x$. For $X\ge1$, since $\psi(x)$ is the sum of the $\Lambda(n)$ with $n\le x$,
\[
\begin{aligned}\int_1^X\psi(x)x^{-s-1}\dd x&=\sum_{n\le X}\Lambda(n)\int_n^Xx^{-s-1}\dd x\\&=\frac1s\sum_{n\le X}\frac{\Lambda(n)}{n^s}-\frac{\psi(X)}{sX^s}.\end{aligned}
\]
As $X\to\infty$ the last term is $O(X^{1-\operatorname{Re}s})\to0$, and the sum tends to $-\frac{\zeta'(s)}{\zeta(s)}$ (Lemma~\ref{lem:logderiv}). With $\int_1^\infty x^{-s}\dd x=\frac1{s-1}$,
\[
G(z)=-\frac{\zeta'(s)}{s\,\zeta(s)}-\frac1{s-1}.
\]

\emph{The continuation.} On $V\setminus\{1\}$, $\frac{\zeta'}\zeta=\frac{Z'}Z-\frac1{s-1}$, so
\[
\begin{aligned}G(z)&=-\frac1s\frac{Z'(s)}{Z(s)}+\frac1{s(s-1)}-\frac1{s-1}\eqbr=-\frac1s\Bigl(\frac{Z'(s)}{Z(s)}+1\Bigr).\end{aligned}
\]
The right side, $g(z)$, is holomorphic for $s=z+1$ in $V$, which is an open set containing $\operatorname{Re}z\ge0$. Its value at $z=0$ is $-(Z'(1)+1)=-1-\gamma$.

\emph{Conclusion.} Newman's theorem gives the convergence of the integral, with value $-1-\gamma$, and Lemma~\ref{lem:monotone} with $A=\psi$, $c=1$ gives $\psi(x)\sim x$. Then
$\theta(x)\sim x$, because $\psi(x)-\theta(x)=O(\sqrt x\log x)$ (Proposition~\ref{prop:chebyshev}). Finally, the proof of Corollary~\ref{cor:chebratio} shows that for every $\varepsilon\in\left]0,1\right[$,
\[
\frac{\theta(x)}x\le\frac{\pi(x)\log x}x\le\frac1{1-\varepsilon}\Bigl(\frac{\theta(x)}x+\frac{\log x}{x^\varepsilon}\Bigr),
\]
so $\frac{\pi(x)\log x}x$ has all its limit points in $[1,\frac1{1-\varepsilon}]$ for every $\varepsilon$, hence tends to $1$.
\end{proof}

The value of the integral is a bonus of Newman's theorem: it gives $g(0)$, not only convergence. It is what Chapter~1 needed to show that $R(x)$ converges
(Exercise~\ref{ex:rlimit}). In Gauss's form, $\pi(x)\sim\operatorname{Li}(x)$, since $\operatorname{Li}(x)\sim\frac x{\log x}$ (Exercise~\ref{ex:li}); and the $n$-th prime is about $n\log n$
(Exercise~\ref{ex:nthprime}).

\begin{remark}[how good is the approximation?]
The proof gives no error term. De la Vallée Poussin showed in 1899, with a zero-free region, that $\pi(x)=\operatorname{Li}(x)+O\bigl(xe^{-c\sqrt{\log x}}\bigr)$ for some $c>0$, which is
better than $O\bigl(\frac x{\log^kx}\bigr)$ for every $k$; this is in Chapter~7. Von Koch showed in 1901 that the Riemann hypothesis is equivalent to
$\pi(x)=\operatorname{Li}(x)+O(\sqrt x\log x)$. The best unconditional error, from Vinogradov and Korobov (1958), is still of the form $O\bigl(xe^{-c(\log x)^{3/5}(\log\log x)^{-1/5}}\bigr)$.
In Figure~\ref{fig:pnt}, $\operatorname{Li}(x)>\pi(x)$ throughout; Littlewood proved in 1914 that the difference nevertheless changes sign infinitely often, but no sign change is known
explicitly, and the first one is believed to occur near $1.4\cdot10^{316}$.
\end{remark}

\begin{remark}[elementary proofs]
In 1948 Selberg and Erdős found proofs that use no complex analysis, starting from Selberg's identity $\sum_{p\le x}\log^2p+\sum_{pq\le x}\log p\log q=2x\log x+O(x)$. They are
``elementary'' only in that sense: they are longer and harder than Newman's. See Tenenbaum's book (Chapter~I.4) or Apostol's (Chapter~4 discusses the equivalent forms used there).
\end{remark}

\section{The mean of the Möbius function}

The same argument gives a statement that Euler had guessed: $\sum\frac{\mu(n)}n=0$. Chapter~1 (Exercise~\ref{ex:mobiusmean}) showed that the partial sums are bounded by $1$. The extra
ingredient replaces monotonicity.

\begin{lemma}\label{lem:slowosc}
Let $F\colon[0,\infty[\to\R$ be such that $\int_0^\infty F(t)\dd t$ converges, and suppose that $|F(u)-F(t)|\le2(u-t)+e^{-t}$ whenever $t\le u\le t+1$. Then $F(t)\to0$.
\end{lemma}
\begin{proof}
Let $\varepsilon\in\left]0,1\right]$, and suppose that $|F(t)|\ge\varepsilon$ for arbitrarily large $t$; say $F(t)\ge\varepsilon$ (the case $F(t)\le-\varepsilon$ is symmetric), with $e^{-t}\le\frac\varepsilon4$. For
$u\in[t,t+\frac\varepsilon4]$, $F(u)\ge\varepsilon-\frac\varepsilon2-\frac\varepsilon4=\frac\varepsilon4$, so $\int_t^{t+\varepsilon/4}F\ge\frac{\varepsilon^2}{16}$. This happens for arbitrarily large $t$, which contradicts the Cauchy criterion.
\end{proof}

\begin{theorem}[von Mangoldt, 1897]\label{thm:mobiussum}
The series $\sum_{n\ge1}\frac{\mu(n)}n$ converges, and its sum is $0$.
\end{theorem}
\begin{proof}
Let $A(x)=\sum_{n\le x}\frac{\mu(n)}n$, so $|A|\le1$, and $F(t)=A(e^t)$, a bounded step function. For $\operatorname{Re}z>0$ and $X\ge1$, as in the last proof,
\[
\begin{aligned}\int_1^XA(x)x^{-z-1}\dd x&=\sum_{n\le X}\frac{\mu(n)}n\int_n^Xx^{-z-1}\dd x\\&=\frac1z\sum_{n\le X}\frac{\mu(n)}{n^{z+1}}-\frac{A(X)}{zX^z},\end{aligned}
\]
and as $X\to\infty$ this tends to $\frac1{z\,\zeta(z+1)}$ (Volume~5, Corollary~12.4). So the Laplace transform of $F$ is $G(z)=\frac1{Z(z+1)}$, which is holomorphic for $z+1\in V$ and equals $1$
at $z=0$. By Newman's theorem, $\int_0^\infty F(t)\dd t=\int_1^\infty\frac{A(x)}x\dd x$ converges (to $1$).

Now let $t\le u\le t+1$. Then $|F(u)-F(t)|\le\sum_{e^t<n\le e^u}\frac1n$, a sum of at most $e^u-e^t+1$ terms, each less than $e^{-t}$; so it is at most
$e^{u-t}-1+e^{-t}\le2(u-t)+e^{-t}$, using $e^v-1\le2v$ for $0\le v\le1$ (convexity: $e^v-1\le(e-1)v$). Lemma~\ref{lem:slowosc} gives $F(t)\to0$, that is, $A(x)\to0$.
\end{proof}

Numerically, $A(10^4)=-0.0021$, $A(10^6)=0.00020$, $A(10^7)=0.00010$. The convergence is slow and the sign wanders, as for $M(x)=\sum_{n\le x}\mu(n)$, which is $o(x)$
(Exercise~\ref{ex:mertensfn}).

\begin{remark}
Theorems~\ref{thm:pnt} and~\ref{thm:mobiussum} are ``equivalent'': each can be deduced from the other, and both from $M(x)=o(x)$, by elementary manipulations of Dirichlet
convolutions (Landau). This is why number theorists call a statement \emph{equivalent to the prime number theorem} when it is as hard as the theorem, no harder and no easier.
Apostol's Chapter~4 proves several of these equivalences. The Riemann hypothesis is equivalent to $M(x)=O(x^{1/2+\varepsilon})$ for every $\varepsilon>0$.
\end{remark}

\section{Primes in arithmetic progressions}

The method needs only three things: a bounded function, a transform continuing across the line, and monotonicity. For a class $a$ prime to $q$, let
\[
\begin{gathered}\psi(x;q,a)=\sum_{\substack{n\le x\\n\equiv a\,(q)}}\Lambda(n),\qq\pi(x;q,a)=\#\{p\le x:p\equiv a\ (q)\}.\end{gathered}
\]
They increase, and $\psi(x;q,a)\le\psi(x)=O(x)$. Theorem~\ref{thm:line} provides the continuation.

\begin{theorem}[the prime number theorem for arithmetic progressions]\label{thm:pntap}
Let $q\ge1$ and $a$ prime to $q$.
\begin{enumerate}[label=(\alph*)]
\item $\psi(x;q,a)\sim\frac x{\varphi(q)}$ and $\pi(x;q,a)\sim\frac1{\varphi(q)}\frac x{\log x}$.
\item As $x\to\infty$, both
\[
\begin{gathered}\sum_{\substack{n\le x\\n\equiv a\,(q)}}\frac{\Lambda(n)}n-\frac{\log x}{\varphi(q)}\quad\text{and}\qq\sum_{\substack{p\le x\\p\equiv a\,(q)}}\frac{\log p}p-\frac{\log x}{\varphi(q)}\end{gathered}
\]
have limits. In particular the function $R_a$ of Chapter~4 converges.
\end{enumerate}
\end{theorem}
\begin{proof}
(a) Let $f(t)=\psi(e^t;q,a)e^{-t}-\frac1{\varphi(q)}$, bounded and piecewise continuous. Exactly as for $\psi$, for $s=z+1$ with $\operatorname{Re}z>0$, its Laplace transform is
\[
G(z)=\frac1s\sum_{n\equiv a\,(q)}\frac{\Lambda(n)}{n^s}-\frac1{\varphi(q)(s-1)}.
\]
Filtering the class with characters (Proposition~\ref{prop:orth}(b)) and Lemma~\ref{lem:logderiv} give
\[
\sum_{n\equiv a\,(q)}\frac{\Lambda(n)}{n^s}=\frac1{\varphi(q)}\sum_\chi\overline{\chi(a)}\Bigl(-\frac{L'(s,\chi)}{L(s,\chi)}\Bigr).
\]
Let $Z_q(s)=(s-1)L(s,\chi_0)$, holomorphic on $\operatorname{Re}s>0$ with $Z_q(1)=\frac{\varphi(q)}q$ (Theorem~\ref{thm:lbasic}(c)). By Theorem~\ref{thm:line}, $Z_q$ and the $L(s,\chi)$, $\chi\ne\chi_0$, have no
zero on $\operatorname{Re}s\ge1$; by continuity they have none on an open set $V_q$ containing that half-plane. On $V_q\setminus\{1\}$, $-\frac{L'}L(s,\chi_0)=-\frac{Z_q'}{Z_q}(s)+\frac1{s-1}$, and the
term $\frac1{s-1}$, which comes with $\overline{\chi_0(a)}=1$, combines with the last term of $G$ as before:
\[
\begin{aligned}G(z)&=\frac1{\varphi(q)s}\Bigl(-\frac{Z_q'(s)}{Z_q(s)}-\sum_{\chi\ne\chi_0}\overline{\chi(a)}\frac{L'(s,\chi)}{L(s,\chi)}\Bigr)\eqbr-\frac1{\varphi(q)s},\end{aligned}
\]
which is holomorphic for $s$ in $V_q$. Newman's theorem shows that $\int_1^\infty\bigl(\psi(x;q,a)-\frac x{\varphi(q)}\bigr)\frac{\mathrm dx}{x^2}$ converges, and Lemma~\ref{lem:monotone} with
$c=\frac1{\varphi(q)}$ gives $\psi(x;q,a)\sim\frac x{\varphi(q)}$. The passage to $\pi(x;q,a)$ is word for word that of Theorem~\ref{thm:pnt}, restricted to the primes of the class.

(b) Write $I_a$ for the integral above. By Lemma~\ref{lem:abelint} (with $y=1$, where the sum is $0$),
\[
\begin{aligned}\sum_{\substack{n\le x\\n\equiv a}}\frac{\Lambda(n)}n&=\frac{\psi(x;q,a)}x+\int_1^x\frac{\psi(t;q,a)}{t^2}\dd t\\&=\frac{\psi(x;q,a)}x+\frac{\log x}{\varphi(q)}\eqbr+\int_1^x\frac{\psi(t;q,a)-t/\varphi(q)}{t^2}\dd t,\end{aligned}
\]
so the first difference tends to $\frac1{\varphi(q)}+I_a$ by (a). The second differs from the first by $\sum\frac{\log p}{p^k}$ over the $p^k\le x$ with $k\ge2$ and $p^k\equiv a$, a sum of positive
terms bounded by the convergent series $\sum_p\frac{\log p}{p(p-1)}$; so it converges too.
\end{proof}

So the classes are equally rich in the plain counting sense, not only in the logarithmic sense of Chapter~4. Exercise~\ref{ex:lprimeseries} identifies the limits in terms of
$\frac{L'}L(1,\chi)$; for $q=3$ this explains the gap of Figure~\ref{fig:mertensprog}. The \emph{order} in which the classes come is another matter, which Chapter~7 takes up.

\section{Exercises}

\begin{exercise}[$\star$]\label{ex:li}
Show that $\operatorname{Li}(x)=\frac x{\log x}+\frac x{\log^2x}+O\bigl(\frac x{\log^3x}\bigr)$. Deduce that $\pi(x)\sim\operatorname{Li}(x)$, and explain the gap between the two curves of Figure~\ref{fig:pnt},
assuming that $\pi(x)-\operatorname{Li}(x)$ is much smaller than $\frac x{\log^2x}$.
\end{exercise}

\begin{exercise}[$\star\star$]\label{ex:rlimit}
Show that $\sum_{n\le x}\frac{\Lambda(n)}n=\log x-\gamma+o(1)$. Deduce that the function $R(x)=\sum_{p\le x}\frac{\log p}p-\log x$ of Chapter~1 tends to
$-\gamma-\sum_p\frac{\log p}{p(p-1)}=-1.33258\ldots$
\end{exercise}

\begin{exercise}[$\star$]\label{ex:mertensfn}
Deduce from Theorem~\ref{thm:mobiussum} that $M(x)=\sum_{n\le x}\mu(n)=o(x)$. In words: the squarefree integers with an even number of prime factors are as frequent as those with
an odd number.
\end{exercise}

\begin{exercise}[$\star\star$, the signature exercise]\label{ex:nthprime}
Let $p_n$ be the $n$-th prime. Show that $p_n\sim n\log n$. (At $n=10^6$, $p_n=15\,485\,863$ and $\frac{p_n}{n\log n}=1.121$; Cipolla's refinement $n(\log n+\log\log n-1)$ is off by $0.29\%$.)
\end{exercise}

\begin{exercise}[$\star$]\label{ex:shortint}
Let $\varepsilon>0$. Show that the number of primes in $]x,(1+\varepsilon)x]$ is asymptotic to $\frac{\varepsilon x}{\log x}$. Compare with Bertrand's postulate (Volume~2, Theorem~2.4).
\end{exercise}

\begin{exercise}[$\star$]\label{ex:lcm}
Show that $\operatorname{lcm}(1,2,\dots,n)^{1/n}\to e$. (The values for $n=100$, $1000$ and $10^8$ are $2.561$, $2.709$ and $2.71823$.)
\end{exercise}

\begin{exercise}[$\star$]\label{ex:ratios}
Show that the quotients $\frac pq$ of two primes are dense in $]0,\infty[$.
\end{exercise}

\begin{exercise}[$\star\star$]\label{ex:lprimeseries}
Let $\chi\ne\chi_0$ be a character modulo $q$. Show that
\[
\sum_{n\ge1}\frac{\chi(n)\Lambda(n)}n=-\frac{L'(1,\chi)}{L(1,\chi)},
\]
the series being convergent. For $q=4$ the value is $-0.24561$; check that it equals $\lim\bigl(\sum_{n\le x,\,n\equiv1}\frac{\Lambda(n)}n-\sum_{n\le x,\,n\equiv3}\frac{\Lambda(n)}n\bigr)$.
\end{exercise}

\begin{exercise}[$\star\star\star$]\label{ex:semiprimes}
Let $\pi_2(x)$ be the number of integers $n\le x$ that are products of exactly two primes, $n=pq$ with $p\le q$ (Landau, 1900). Show that
\[
\pi_2(x)\sim\frac{x\log\log x}{\log x}.
\]
(The ratio is $1.102$ at $x=10^8$: the convergence is very slow.)
\end{exercise}

\begin{exercise}[$\star\star$]\label{pb:gaussprimes}
Let $G(x)$ be the number of Gaussian primes $\pi\in\Z[i]$ with $N(\pi)\le x$, all associates counted.
\begin{enumerate}[label=(\alph*)]
\item Show that $G(x)=4\bigl(1+2\pi(x;4,1)+\pi(\sqrt x;4,3)\bigr)$, where $\pi(y;4,a)$ counts the primes $p\le y$ with
$p\equiv a\pmod4$.
\item Show that $G(x)\sim\frac{4x}{\log x}$. Since about $\pi x$ Gaussian integers have norm at most $x$, what is the
``probability'' that a Gaussian integer of norm near $N$ is prime?
\item Show that $\sum\frac1{N(\pi)}=\log\log x+O(1)$, the sum over the Gaussian primes of norm at most $x$, one in each
class of associates.
\end{enumerate}
\end{exercise}

\begin{exercise}[$\star$]\label{pb:findprime}
To make an RSA key, one draws random integers in $[N,2N]$ until one is prime.
\begin{enumerate}[label=(\alph*)]
\item Show that the expected number of draws is asymptotic to $\log N$.
\item If one only draws integers prime to $30$, show that it becomes $\sim\frac4{15}\log N$. With only integers prime to all
primes $\le y$ ($y$ fixed, then large), what is the saving?
\item Estimate both numbers for $N=2^{1023}$.
\end{enumerate}
\end{exercise}

\begin{exercise}[$\star\star$]\label{pb:sidon}
\emph{Sidon sets.} A set $A$ of integers is a \emph{Sidon set} if all the sums $a+b$, $a\le b$ in $A$, are distinct. Let $F(n)$
be the largest size of a Sidon set in $\{1,\dots,n\}$.
\begin{enumerate}[label=(\alph*)]
\item Show that $F(n)<\sqrt{2n}+1$.
\item (Erdős and Turán, 1941.) Let $p$ be an odd prime. Show that $\{2pk+(k^2\bmod p):0\le k<p\}$ is a Sidon set, with all elements
below $2p^2$.
\item (Erdős and Turán.) Show that $F(n)\le\sqrt n+O(n^{1/4})$. Count the pairs of elements of $A$ that fall in a common window
of $m$ consecutive integers, in two ways.
\item Using Singer sets and the prime number theorem, show that $F(n)\sim\sqrt n$.
\end{enumerate}
\end{exercise}

\begin{exercise}[$\star\star$]\label{pb:nsq1sqfree}
\emph{Squarefree values of $n^2+1$.}
\begin{enumerate}[label=(\alph*)]
\item Show that the density of the $n$ such that $n^2+1$ is squarefree is
\[
\prod_{p\equiv1\ (4)}\Bigl(1-\frac2{p^2}\Bigr)=0.8948\ldots
\]
\item Why is this so much larger than the density $\frac6{\pi^2}=0.6079\ldots$ of all squarefree integers?
\end{enumerate}
\end{exercise}

\section{Hints}

\begin{hint}{ex:li}
Integrate by parts twice. For the remainder $\int_2^x\frac{\mathrm dt}{\log^3t}$, split at $\sqrt x$.
\end{hint}
\begin{hint}{ex:rlimit}
Write $\sum_{n\le x}\frac{\Lambda(n)}n$ with Abel summation in terms of $\psi$, and use the value of the integral in Theorem~\ref{thm:pnt}.
\end{hint}
\begin{hint}{ex:mertensfn}
$\mu(n)=n\cdot\frac{\mu(n)}n$: apply Abel summation with $A(x)=\sum_{n\le x}\frac{\mu(n)}n$ and $f(t)=t$.
\end{hint}
\begin{hint}{ex:nthprime}
$\pi(p_n)=n$. Take logarithms in $n\sim\frac{p_n}{\log p_n}$ to compare $\log n$ and $\log p_n$.
\end{hint}
\begin{hint}{ex:shortint}
Compute $\pi((1+\varepsilon)x)$ and $\pi(x)$ with the prime number theorem, and note that $\log((1+\varepsilon)x)\sim\log x$.
\end{hint}
\begin{hint}{ex:lcm}
$\log\operatorname{lcm}(1,\dots,n)=\psi(n)$.
\end{hint}
\begin{hint}{ex:ratios}
Given $0<a<b$, look for a prime $p$ in $]aq,bq[$ for a large prime $q$.
\end{hint}
\begin{hint}{ex:lprimeseries}
Apply Newman's theorem to $f(t)=e^{-t}\sum_{n\le e^t}\chi(n)\Lambda(n)$, which is bounded and complex-valued, and show that $\sum_{n\le x}\chi(n)\Lambda(n)=o(x)$ with Theorem~\ref{thm:pntap}.
\end{hint}
\begin{hint}{ex:semiprimes}
$\pi_2(x)=\sum_{p\le\sqrt x}\bigl(\pi(\frac xp)-\pi(p)+1\bigr)$. For $p\le\sqrt x$, $\frac xp\ge\sqrt x$ is large. Then estimate $\sum_{p\le\sqrt x}\frac1{p\log(x/p)}$, splitting at $p=x^\varepsilon$.
\end{hint}

\begin{hint}{pb:gaussprimes}
(a) The Gaussian primes up to units are $1+i$, the $\pi$ and $\bar\pi$ above each $p\equiv1\pmod4$, and the $p\equiv3\pmod4$. (b)
and (c): the prime number theorem and Dirichlet's theorem in progressions.
\end{hint}

\begin{hint}{pb:findprime}
A geometric number of draws. (b) Mertens's product $\prod_{p\le y}(1-\frac1p)\sim\frac{e^{-\gamma}}{\log y}$.
\end{hint}

\begin{hint}{pb:sidon}
(b) Equal sums give equal $k+l$ and equal $k^2+l^2$ modulo $p$. (c) Window length $m\approx n^{3/4}$; a difference $d<m$ occurs at
most once.
\end{hint}

\begin{hint}{pb:nsq1sqfree}
Count the roots of $x^2\equiv-1$ modulo $p^2$. For the tail, $p^2\mid n^2+1$ with $n\le N$ forces $p\le N$.
\end{hint}

\section{Solutions}

\begin{solution}{ex:li}
\idea Each integration by parts trades $\frac1{\log^kt}$ for $\frac1{\log^{k+1}t}$, which is smaller; the last integral is small because $\log t$ is comparable to $\log x$ on most of $[2,x]$.
\writeup With $u=\frac1{\log^kt}$, $u'=-\frac k{t\log^{k+1}t}$:
\[
\int_2^x\frac{\mathrm dt}{\log^kt}=\frac x{\log^kx}-\frac2{\log^k2}+k\int_2^x\frac{\mathrm dt}{\log^{k+1}t}.
\]
With $k=1$ and then $k=2$, $\operatorname{Li}(x)=\frac x{\log x}+\frac x{\log^2x}+2\int_2^x\frac{\mathrm dt}{\log^3t}+O(1)$. On $[2,\sqrt x]$ the integrand is at most $\frac1{\log^32}$, and on $[\sqrt x,x]$ at most
$\frac8{\log^3x}$; so the integral is at most $\frac{\sqrt x}{\log^32}+\frac{8x}{\log^3x}=O\bigl(\frac x{\log^3x}\bigr)$. In particular $\operatorname{Li}(x)\sim\frac x{\log x}\sim\pi(x)$.

If $\pi(x)=\operatorname{Li}(x)+o\bigl(\frac x{\log^2x}\bigr)$, then $\frac{\pi(x)\log x}x=1+\frac1{\log x}+O\bigl(\frac1{\log^2x}\bigr)$: the coral curve approaches $1$ only like $\frac1{\log x}$, which is still
$0.054$ at $x=10^8$. (The observed excess, $0.061$, includes the next term, $\frac2{\log^2x}$.) The hypothesis is true: by Chapter~7, the error is much smaller than any power of $\frac1{\log x}$ times $x$.
\end{solution}

\begin{solution}{ex:rlimit}
\idea Abel summation turns the sum into $\frac{\psi(x)}x$ plus an integral of $\frac{\psi(t)}{t^2}$; subtracting $\log x=\int_1^x\frac{\mathrm dt}t$ leaves the integral of Theorem~\ref{thm:pnt}, whose value we know.
\writeup By Lemma~\ref{lem:abelint} with $a_n=\Lambda(n)$, $f(t)=\frac1t$ and $y=1$,
\[
\begin{aligned}\sum_{n\le x}\frac{\Lambda(n)}n&=\frac{\psi(x)}x+\int_1^x\frac{\psi(t)}{t^2}\dd t\\&=\frac{\psi(x)}x+\log x+\int_1^x\frac{\psi(t)-t}{t^2}\dd t.\end{aligned}
\]
By Theorem~\ref{thm:pnt}, $\frac{\psi(x)}x\to1$ and the integral tends to $-1-\gamma$. So the sum is $\log x+1-1-\gamma+o(1)=\log x-\gamma+o(1)$.

The sum over primes differs from this by the sum of $\frac{\log p}{p^k}$ over the $p^k\le x$ with $k\ge2$, which increases to
\[
\sum_p\log p\sum_{k\ge2}p^{-k}=\sum_p\frac{\log p}{p(p-1)}=0.75536\ldots
\]
So
$R(x)\to-\gamma-0.75536\ldots=-1.33258\ldots$, as announced in Chapter~1. At $x=10^8$, $\sum_{n\le x}\frac{\Lambda(n)}n-\log x=-0.577234$, against $-\gamma=-0.577216$.
\end{solution}

\begin{solution}{ex:mertensfn}
\idea Partial summation writes $M(x)$ as $xA(x)$ minus an integral of $A$, and both are $o(x)$ when $A\to0$.
\writeup By Lemma~\ref{lem:abelint} with $a_n=\frac{\mu(n)}n$, $f(t)=t$ and $y=1$: $M(x)=xA(x)-\int_1^xA(t)\dd t$. Let $\varepsilon>0$ and $t_0$ such that $|A(t)|\le\varepsilon$ for $t\ge t_0$
(Theorem~\ref{thm:mobiussum}). Since $|A|\le1$, $|M(x)|\le\varepsilon x+t_0+\varepsilon x$ for $x\ge t_0$, so $\limsup\frac{|M(x)|}x\le2\varepsilon$. For squarefree $n$, $\mu(n)=(-1)^{\omega(n)}$, so $M(x)$
is the number of squarefree $n\le x$ with $\omega(n)$ even minus the number with $\omega(n)$ odd; each count is $\frac3{\pi^2}x+o(x)$ (Chapter~2 for their sum).
\end{solution}

\begin{solution}{ex:nthprime}
\idea Invert $\pi(x)\sim\frac x{\log x}$: the inverse function of $\frac x{\log x}$ is about $n\log n$, because taking logarithms kills the $\log$ factor.
\writeup Since $\pi(p_n)=n$, the prime number theorem gives $\frac{n\log p_n}{p_n}\to1$. Taking logarithms, $\log n+\log\log p_n-\log p_n\to0$, so
$\frac{\log n}{\log p_n}=1-\frac{\log\log p_n}{\log p_n}+o\bigl(\frac1{\log p_n}\bigr)\to1$. Hence $p_n\sim n\log p_n\sim n\log n$. The slow convergence (ratio $1.121$ at $n=10^6$) comes from the
neglected $\log\log p_n$: pushing the same computation one step further gives $p_n=n(\log n+\log\log n-1+o(1))$, Cipolla's formula (1902).
\end{solution}

\begin{solution}{ex:shortint}
\idea The density $\frac1{\log t}$ hardly changes on $[x,(1+\varepsilon)x]$.
\writeup $\pi((1+\varepsilon)x)=\frac{(1+\varepsilon)x}{\log x+\log(1+\varepsilon)}(1+o(1))=\frac{(1+\varepsilon)x}{\log x}(1+o(1))$, and $\pi(x)=\frac x{\log x}(1+o(1))$. The difference is
$\frac{\varepsilon x}{\log x}+o\bigl(\frac x{\log x}\bigr)\sim\frac{\varepsilon x}{\log x}$. In particular for every $\varepsilon>0$ there is a prime between $x$ and $(1+\varepsilon)x$ for $x$ large enough, while Bertrand's postulate
gives one between $n$ and $2n$ for every $n\ge1$: it is weaker for large $n$, but explicit.
\end{solution}

\begin{solution}{ex:lcm}
\idea The exponent of $p$ in $\operatorname{lcm}(1,\dots,n)$ is the largest $k$ with $p^k\le n$.
\writeup So $\log\operatorname{lcm}(1,\dots,n)=\sum_{p^k\le n}\log p=\psi(n)$, each prime $p$ counted once for each power $p^k\le n$. Then $\frac1n\log\operatorname{lcm}(1,\dots,n)=\frac{\psi(n)}n\to1$. Chebyshev's
bounds (Chapter~1) only give that this lies between two constants; the limit $1$ is exactly the prime number theorem in the form $\psi(x)\sim x$.
\end{solution}

\begin{solution}{ex:ratios}
\idea A prime $q$ times a short interval of ratios $]a,b[$ is a long interval, so it contains a prime $p$.
\writeup Let $0<a<b$, and $\varepsilon>0$ with $(1+\varepsilon)a<b$. By Exercise~\ref{ex:shortint}, there is $x_0$ such that $]x,(1+\varepsilon)x]$ contains a prime for $x\ge x_0$. Take a prime $q$ with
$aq\ge x_0$, and a prime $p\in]aq,(1+\varepsilon)aq]$. Then $a<\frac pq\le(1+\varepsilon)a<b$.
\end{solution}

\begin{solution}{ex:lprimeseries}
\idea This is Theorem~\ref{thm:pnt}'s computation of $g(0)$ with $\chi$ inserted: since $L(s,\chi)$ has no pole, there is no main term, and the integral is the series.
\writeup Let $\psi(x,\chi)=\sum_{n\le x}\chi(n)\Lambda(n)$. Filtering by classes,
\[
\psi(x,\chi)=\sum_a\chi(a)\psi(x;q,a),
\]
summed over the classes prime to $q$; by Theorem~\ref{thm:pntap}(a), it is
$\frac x{\varphi(q)}\sum_a\chi(a)+o(x)=o(x)$, by Proposition~\ref{prop:orth}(a). Let $f(t)=e^{-t}\psi(e^t,\chi)$, bounded by $e^{-t}\psi(e^t)$ and piecewise continuous. As in Theorem~\ref{thm:pnt},
its Laplace transform is $G(z)=-\frac1s\frac{L'(s,\chi)}{L(s,\chi)}$ with $s=z+1$, and it continues holomorphically to a neighbourhood of $\operatorname{Re}z\ge0$ because $L(s,\chi)$ is holomorphic on
$\operatorname{Re}s>0$ and does not vanish on $\operatorname{Re}s\ge1$ (Theorem~\ref{thm:line}). Newman's theorem gives $\int_1^\infty\frac{\psi(x,\chi)}{x^2}\dd x=-\frac{L'(1,\chi)}{L(1,\chi)}$. By Abel summation,
$\sum_{n\le x}\frac{\chi(n)\Lambda(n)}n=\frac{\psi(x,\chi)}x+\int_1^x\frac{\psi(t,\chi)}{t^2}\dd t$, which tends to $0-\frac{L'(1,\chi)}{L(1,\chi)}$.

For $q=4$, $\chi(n)=1,-1$ on the classes $1,3$, so the series is the limit of the difference in the statement. Numerically $\frac{L'}L(1,\chi_{-4})=0.24561$ (a classical closed form is
$\gamma+2\log2+3\log\pi-4\log\Gamma(\frac14)$), and the partial sum up to $10^8$ is $-0.24558$. For $q=3$ the same computation gives $\frac{L'}L(1,\chi_{-3})=0.36828$, the part of the gap of
Figure~\ref{fig:mertensprog} that does not come from prime powers.
\end{solution}

\begin{solution}{ex:semiprimes}
\idea Count by the smaller prime $p\le\sqrt x$: the larger one is any prime in $[p,\frac xp]$. The prime number theorem applies uniformly, since $\frac xp\ge\sqrt x$, and the resulting sum
over $p$ is dominated by small $p$, where $\log\frac xp\approx\log x$: that produces $\sum_{p}\frac1p\approx\log\log x$.
\writeup Each $n=pq\le x$ with $p\le q$ has $p\le\sqrt x$ and $q\in[p,\frac xp]$, so $\pi_2(x)=\sum_{p\le\sqrt x}\bigl(\pi(\frac xp)-\pi(p)+1\bigr)$. The terms $\pi(p)-1$ add up to at most
$\pi(\sqrt x)^2=O\bigl(\frac x{\log^2x}\bigr)$, which is negligible.

Write $\pi(y)=(1+\eta(y))\frac y{\log y}$ with $\eta(y)\to0$, and $\eta^*(y)=\sup_{u\ge y}|\eta(u)|\to0$. For $p\le\sqrt x$, $\frac xp\ge\sqrt x$, so
\[
\begin{aligned}\sum_{p\le\sqrt x}\pi\Bigl(\frac xp\Bigr)&=\bigl(1+O(\eta^*(\sqrt x))\bigr)\,x\,S(x),\eqbr S(x)=\sum_{p\le\sqrt x}\frac1{p\log(x/p)}.\end{aligned}
\]
It remains to show that $S(x)\sim\frac{\log\log x}{\log x}$. Since $\log\frac xp\le\log x$, $S(x)\ge\frac1{\log x}\sum_{p\le\sqrt x}\frac1p=\frac{\log\log x+O(1)}{\log x}$ (Theorem~\ref{thm:mertens2}). Conversely, let
$0<\varepsilon<\frac12$. For $p\le x^\varepsilon$, $\log\frac xp\ge(1-\varepsilon)\log x$, and these terms add up to at most $\frac{\log\log x+O(1)}{(1-\varepsilon)\log x}$. For $x^\varepsilon<p\le\sqrt x$,
$\log\frac xp\ge\frac12\log x$, and by Theorem~\ref{thm:mertens2} $\sum_{x^\varepsilon<p\le\sqrt x}\frac1p=\log\frac1{2\varepsilon}+o(1)$, so these terms add up to $O_\varepsilon\bigl(\frac1{\log x}\bigr)$. Hence
$\limsup\frac{S(x)\log x}{\log\log x}\le\frac1{1-\varepsilon}$ for every $\varepsilon$, and $S(x)\sim\frac{\log\log x}{\log x}$.

The ratio $1.102$ at $10^8$ reflects the neglected terms of relative size $\frac1{\log\log x}$, which decay extremely slowly. More generally, Landau showed that the number of $n\le x$ with
exactly $k$ prime factors is $\sim\frac x{\log x}\frac{(\log\log x)^{k-1}}{(k-1)!}$: a Poisson law in disguise, which Chapter~9 explains.
\end{solution}

\begin{solution}{pb:gaussprimes}
\idea Volume~2 lists the Gaussian primes: $1+i$ (norm $2$), a conjugate pair of norm $p$ for each $p\equiv1\pmod4$, and the
primes $q\equiv3\pmod4$ themselves (norm $q^2$), each with four associates. Counting by norm reduces to counting primes in
progressions, and the asymptotics come from Volume~6.

\writeup (a) By Volume~2, Chapter~10, every Gaussian prime is associate to exactly one of: $1+i$; $\pi$ or $\bar\pi$, where
$p=\pi\bar\pi\equiv1\pmod4$ (and $\pi,\bar\pi$ are not associate); $q\equiv3\pmod4$ prime in $\Z$. Their norms are $2$, $p$, $q^2$, and
each class has $4$ elements (the units $\pm1,\pm i$). So $G(x)=4\bigl(1+2\pi(x;4,1)+\pi(\sqrt x;4,3)\bigr)$ for $x\ge2$.

(b) By the prime number theorem in progressions (Volume~6, Chapter~5), $\pi(x;4,1)\sim\frac x{2\log x}$, and
$\pi(\sqrt x;4,3)\le\sqrt x$. So $G(x)\sim\frac{4x}{\log x}$. The number of Gaussian integers of norm at most $x$ is the number of lattice
points in the disc of radius $\sqrt x$, $\sim\pi x$ (Volume~2, Chapter~10, or a direct count). By the same differencing as for
$\pi(x)$, the density of Gaussian primes near norm $N$ is $\frac4{\pi\log N}$: in the plane, a Gaussian integer of norm near $N$
is prime with probability about $\frac{4}{\pi}\cdot\frac1{\log N}$.

(c) The sum is $\frac12+2\sum_{p\le x,p\equiv1(4)}\frac1p+\sum_{q\equiv3(4),q^2\le x}\frac1{q^2}$. The last sum converges, and by
Dirichlet's theorem in Mertens form (Volume~6, Chapter~4, or by partial summation from the
$\frac{\log p}p$ version), $\sum_{p\le x,p\equiv1(4)}\frac1p=\frac12\log\log x+O(1)$. The total is $\log\log x+O(1)$.
\tools Volume~2, Chapter~10 (Gaussian primes); Volume~6, Chapters~4 and~5 (Dirichlet, prime number theorem in progressions).
\leads (c) is Mertens's theorem for the ring $\Z[i]$, and the analogue holds for every number field (Landau's prime ideal
theorem, Volume~13). The Gaussian primes are equidistributed in angle (Hecke, 1920), and Duke, Friedlander and Iwaniec showed that
there are Gaussian primes in very thin sectors. Whether one can walk to infinity on Gaussian primes with bounded steps (the
\emph{Gaussian moat problem}) is open.
\end{solution}

\begin{solution}{pb:findprime}
\idea A draw succeeds with probability $\frac{\pi(2N)-\pi(N)}{N+1}$, and the number of draws until the first success is geometric, with
mean the inverse of that. The prime number theorem gives the probability; sieving out small factors raises it by the factor by
which they thin out the candidates.

\writeup (a) By the prime number theorem (Volume~6, Chapter~5), $\pi(2N)-\pi(N)\sim\frac{2N}{\log2N}-\frac N{\log N}\sim\frac N{\log N}$. The
number of independent uniform draws until a prime is geometric with parameter $p_N\sim\frac1{\log N}$, of mean $\frac1{p_N}\sim\log N$
(Volume~4, Chapter~2).

(b) The integers of $[N,2N]$ prime to $30$ number $\frac8{30}N+O(1)$, and all primes $>5$ are among them. So the success probability
becomes $\sim\frac{N/\log N}{(4/15)N}$, and the mean number of draws $\sim\frac4{15}\log N$. With all primes $\le y$, the factor is
$\prod_{p\le y}(1-\frac1p)$, which by Mertens's theorem (Volume~6, Chapter~1) is $\sim\frac{e^{-\gamma}}{\log y}$: sieving to $y$ divides the
work by about $1.78\log y$.

(c) $\log2^{1023}=1023\log2\approx709$, so about $709$ draws; with the wheel $30$, about $189$. Each draw that survives costs one
Miller--Rabin test, the others a few trial divisions.
\tools Volume~6, Chapters~1 and~5 (Mertens, prime number theorem); Volume~4, Chapter~2 (geometric law).
\leads Real implementations sieve an interval rather than draw independently, which biases the choice: primes after a long gap are
chosen more often. Brandt and Damgård (1993) and later Fouque and Tibouchi analysed this bias and its (small) effect on security.
\end{solution}

\begin{solution}{pb:sidon}
\idea A Sidon set is a set with distinct differences. There are only $n-1$ possible differences, which gives $\sqrt{2n}$; counting
only short differences, in windows, saves the factor $\sqrt2$. Singer sets reach the bound.

\writeup (a) $a+b=a'+b'$ is the same as $a-a'=b'-b$, so in a Sidon set the $\binom k2$ positive differences are distinct, and they
lie in $\{1,\dots,n-1\}$. So $k(k-1)\le2(n-1)$ and $k<\sqrt{2n}+1$.

(b) Let $a_k=2pk+r_k$ with $r_k=k^2\bmod p\in[0,p)$. Then $a_k+a_l=2p(k+l)+(r_k+r_l)$ with $0\le r_k+r_l<2p$, so $a_k+a_l=a_{k'}+a_{l'}$
forces $k+l=k'+l'$ and $r_k+r_l=r_{k'}+r_{l'}$, hence $k^2+l^2\equiv k'^2+l'^2\pmod p$. With $s=k+l$, $2kl=s^2-(k^2+l^2)$, so
$kl\equiv k'l'$, and $\{k,l\}$, $\{k',l'\}$ are both the roots of $X^2-sX+kl$ in $\mathbb F_p$; as $0\le k,l,k',l'<p$, they are the
same pair. The largest element is below $2p(p-1)+p<2p^2$.

(c) Let $A\subset\{1,\dots,n\}$ be Sidon, $|A|=k$, and $1\le m\le n$. For $j=-m+1,\dots,n-1$ (that is $n+m-1$ values) let
$I_j=\{j+1,\dots,j+m\}$ and $A_j=|A\cap I_j|$. Each $a\in A$ lies in exactly $m$ windows, so $\sum_jA_j=km$, and by Cauchy--Schwarz
\[
\sum_j\binom{A_j}2\ge\frac12\Bigl(\frac{(km)^2}{n+m-1}-km\Bigr).
\]
On the other hand $\sum_j\binom{A_j}2$ counts pairs $a<b$ of $A$ with a common window, with multiplicity: if $b-a=d<m$ they share $m-d$
windows. Each $d$ is the difference of at most one pair, so $\sum_j\binom{A_j}2\le\sum_{d=1}^{m-1}(m-d)=\frac{m(m-1)}2$. Combining,
$\frac{k^2m}{n+m}-k\le m-1$, that is
\[
k^2\le\frac{(n+m)(m+k-1)}m=n\Bigl(1+\frac mn\Bigr)\Bigl(1+\frac{k-1}m\Bigr).
\]
Take $m=\lceil n^{3/4}\rceil$. By (a), $\frac{k-1}m\le\sqrt2\,n^{-1/4}$, so $k^2\le n\bigl(1+O(n^{-1/4})\bigr)$ and
$k\le\sqrt n+O(n^{1/4})$.

(d) Let $q$ be the largest prime with $q^2+q+1\le n$. The Singer set of Problem~\ref{pb:singer}, read as integers in $[0,N)$ and
shifted by $1$, is a Sidon set in $\{1,\dots,n\}$ with $q+1$ elements: distinct differences modulo $N$ are distinct differences
in $\Z$. By the prime number theorem (Volume~6, Chapter~5), consecutive primes satisfy $p_{j+1}/p_j\to1$, so $q\sim\sqrt n$. With
(c), $F(n)\sim\sqrt n$.

\tools Volume~4, Chapter~2 (Cauchy--Schwarz); Problem~\ref{pb:singer}; Volume~6 (the prime number theorem).
\leads Erdős offered \$500 for a proof or disproof of $F(n)=\sqrt n+O(1)$. The best upper bound is
$\sqrt n+0.998n^{1/4}$ for large $n$ (Balogh, Füredi and Roy, 2023), improving the constant $1$ of Lindström (1969). Sidon sets are
also the rulers with all distances distinct, \emph{Golomb rulers}, used to place radio antennas.
\end{solution}

\begin{solution}{pb:nsq1sqfree}
\idea The local factor at $p$ is $1-\rho(p^2)/p^2$, where $\rho(p^2)$ counts the solutions of $n^2\equiv-1\pmod{p^2}$. The novelty
is the tail: the primes up to $N$ are now all relevant, but a polynomial of degree $2$ has few roots, and the number of primes
is $o(N)$.

\writeup Let $\rho(m)$ be the number of $x\bmod m$ with $x^2\equiv-1$. For $p=2$, $x^2+1\equiv1,2\pmod4$, so $\rho(4)=0$. For $p\equiv3\pmod4$,
$-1$ is not a square modulo $p$, so $\rho(p^2)=0$. For $p\equiv1\pmod4$ there are two roots modulo $p$, and each lifts uniquely to
a root modulo $p^2$: if $x_0^2\equiv-1\pmod p$, then $(x_0+tp)^2+1\equiv x_0^2+1+2x_0tp\pmod{p^2}$, and since $2x_0\not\equiv0$ there is
exactly one $t$ modulo $p$. So $\rho(p^2)=2$. In all cases $\rho(p^2)\le2$.

As in Problem~\ref{pb:consecsqfree}, the $n$ with $p^2\nmid n^2+1$ for all $p\le y$ have density $\prod_{p\le y}(1-\rho(p^2)/p^2)$, which is
$\prod_{p\le y,\,p\equiv1(4)}(1-\frac2{p^2})$. For the tail: if $p^2\mid n^2+1$ with $n\le N$, then $p^2\le N^2+1<(N+1)^2$, so $p\le N$. The
number of $n\le N$ with $p^2\mid n^2+1$ is at most $\rho(p^2)(\frac N{p^2}+1)\le\frac{2N}{p^2}+2$. Summing over $y<p\le N$ gives at most
$\frac{2N}y+2\pi(N)$, and $\pi(N)=o(N)$ (Volume~6, Chapter~5, or already Chebyshev's bound). The rest is as before.

(b) The density of squarefree numbers is $\prod_p(1-\frac1{p^2})$: every prime removes a proportion $\frac1{p^2}$, and $p=2$ and $p=3$
remove the most, $\frac14$ and $\frac19$. For $n^2+1$, the primes $2$ and $3$ (and all $p\equiv3\pmod4$) remove nothing: $n^2+1$ is never
divisible by $4$ or by $9$. The first prime that matters is $5$, which removes $\frac2{25}$.

\tools Volume~2, Chapter~6 ($-1$ is a square modulo $p$ exactly when $p\equiv1\pmod4$); Volume~13, Chapter~2 (Hensel's lifting);
Problem~\ref{pb:consecsqfree}.
\leads Estermann (1931) proved (a) with an error term. For a polynomial of degree $3$, squarefree values are known to have the
expected density (Hooley, Erdős for $k$-free values), but for degree $\ge4$ the question is open in general: the tail estimate
fails, because the primes up to $N^{d/2}$ matter and there are too many of them. Granville showed in 1998 that the $abc$
conjecture would settle all degrees.
\end{solution}

\section*{Sources and further reading}

\begin{itemize}
\item D.~J.~Newman, \emph{Simple analytic proof of the prime number theorem}, American Mathematical Monthly 87 (1980), 693--696.
\item D.~Zagier, \emph{Newman's short proof of the prime number theorem}, American Mathematical Monthly 104 (1997), 705--708. The proof of this chapter, in four pages.
\item J.~Hadamard, \emph{Sur la distribution des zéros de la fonction $\zeta(s)$ et ses conséquences arithmétiques}, Bulletin de la Société Mathématique de France 24 (1896), 199--220;
C.-J.~de la Vallée Poussin, \emph{Recherches analytiques sur la théorie des nombres premiers}, Annales de la Société Scientifique de Bruxelles 20 (1896), 183--256.
\item L.~J.~Goldstein, \emph{A history of the prime number theorem}, American Mathematical Monthly 80 (1973), 599--615. Gauss's letter to Encke, in translation, is quoted there.
\item T.~M.~Apostol, \emph{Introduction to Analytic Number Theory}, Springer, 1976, Chapters~4 and~13 (equivalent forms of the theorem; an analytic proof by contour integration).
\item G.~Tenenbaum, \emph{Introduction to Analytic and Probabilistic Number Theory}, 3rd ed., AMS, 2015, Chapters~I.4 and~II.4 (elementary and analytic proofs, error terms).
\item J.~Korevaar, \emph{On Newman's quick way to the prime number theorem}, Mathematical Intelligencer 4 (1982), 108--115. How Newman's theorem fits into Tauberian theory.
\item L.~J.~Goldstein, \emph{A history of the prime number theorem}, American Mathematical Monthly 80 (1973), 599--615, with Gauss's 1849 letter to Encke and his counts up to three million.
\end{itemize}
